\section{Subprogram 4c: Maximality of supersymmetry at the UV fixed point}
\label{sec:subprogram4c}

This subprogram is an optional strengthening of the superconformal
closure obtained in Subprogram~4b. We take that closure as imported
input and ask what additional assumptions would force maximal
supersymmetry at the UV fixed point. The answer is conditional:
under a short list of extra assumptions about non--gravitationality,
interaction, and symmetry selection, the superconformal symmetry must be
\(\mathcal N=4\) (i.e.\ \(N=4\) copies of Weyl supercharges).

\subsection{Inputs and goal}
\label{subsec:4c-inputs-goal}

We begin by collecting the earlier outputs used in this optional
strengthening.

\begin{itemize}
  \item Subprogram~1b produced the Wightman vacuum representation
        \((\mathcal H,\pi,U,\Omega)\) with locality, positivity, and
        positive--energy Poincar\'e covariance
        (Theorem~\ref{thm:subprogram1b-Wightman}).

  \item Subprogram~3c produced the HLS odd charge module on that UV
        boundary representation, together with the operator
        implementation of the odd symmetry derivation
        (Theorem~\ref{thm:Main-HLS} and
        Proposition~\ref{prop:inner}).

  \item Subprogram~4a showed that on the common UV sector, any
        nontrivial odd charge \(Q_\alpha^I\) has dilatation weight
        \(\Delta(Q_\alpha^I)=\tfrac12\)
        (Theorem~\ref{thm:4a-Delta-half}).

  \item Subprogram~4b established the same-\(\mathcal H\)
        superconformal closure of the odd sector: the generators
        \(P_\mu,M_{\mu\nu},D,K_\mu,Q,\bar Q,S,\bar S,R\) close on a
        common invariant core to the full \(4\)D superconformal algebra
        (Theorem~\ref{thm:4d-SConf}). At Main level the same conclusion
        is packaged in Corollary~\ref{cor:Main-postHLS-superconf}.
\end{itemize}

In the present subprogram we:

\begin{enumerate}[(1)]
  \item import from Subprogram~4b the resulting \(4\)D superconformal
        algebra and write \(N\) for its odd multiplicity parameter;

  \item impose that the theory is non–gravitational and has a unique
        stress tensor, together with standard representation–theoretic
        input, to show that \(N\le 4\);

  \item add interaction and nonfactorization assumptions so that \(N\)
        measures the intrinsic supersymmetry of the interacting UV fixed
        point rather than of a theory enlarged by decoupled free
        sectors;

  \item assume a maximality principle and deduce that, within this
        optional strengthening, the emergent supersymmetry is maximal,
        i.e.\ \(N=4\).
\end{enumerate}

Everything beyond the imported superconformal closure of
Subprogram~4b is therefore assumption-driven.

\subsection{Imported superconformal algebra from Subprogram 4b}
\label{subsec:4c-algebraic}

\begin{proposition}[Imported algebraic form of the superconformal symmetry]
\label{prop:4c-su-22N}
Use Theorem~\ref{thm:4d-SConf} as imported input. Then on the same UV
Hilbert space the resulting \(4\)D superconformal algebra is of standard
type: for some integer \(N\ge1\), it is
\(\mathfrak{su}(2,2|N)\) for \(N\neq4\), while in the \(N=4\) case the
standard simple form is obtained after quotienting the central
\(\mathfrak u(1)\), giving \(\mathfrak{psu}(2,2|4)\). Throughout this
section we write \(N\) for this imported multiplicity parameter. When
convenient, we speak informally of the \(N=4\) case as
\(\mathfrak{su}(2,2|4)\), since the central \(\mathfrak u(1)\) does not
affect the arguments below.
\end{proposition}

\begin{proof}
This is exactly the algebraic form isolated in
Theorem~\ref{thm:4d-SConf}; see also
Corollary~\ref{cor:Main-postHLS-superconf} for the corresponding Main
packaging. The Kac--Nahm classification explains why this is the
standard \(4\)D superconformal family, but that classification is not
rederived here.
\end{proof}

At this stage \(N\) is an imported integer label. The purpose of the
remaining subsections is to examine what additional assumptions would
bound or determine it.

\subsection{Stress tensor multiplet and absence of dynamical gravity}
\label{subsec:4c-stresstensor}

We now impose that our theory is purely field–theoretic, with a unique
stress tensor and no dynamical gravity.

\begin{definition}[Stress tensor and its multiplet]
\label{def:4c-T-multiplet}
Let \(T_{\mu\nu}(x)\) be the unique conserved, symmetric, traceless
stress tensor of the boundary Wightman theory constructed in
Subprogram~1b. In the presence of the imported superconformal symmetry
of Proposition~\ref{prop:4c-su-22N}, the \emph{stress tensor multiplet} is the
smallest unitary positive–energy representation of
\(\mathfrak{su}(2,2|N)\) on \(\mathcal H\) whose component fields
include \(T_{\mu\nu}\) and close under the action of all supercharges
\(Q_\alpha^I,\bar Q_{\dot\alpha I}\) and \(S^\alpha{}_I,\bar S^{\dot\alpha I}\).
\end{definition}

The existence and structure of this multiplet is standard: it contains
\(T_{\mu\nu}\), the supercurrents
\(\mathcal J^\mu{}_\alpha{}^I,\bar{\mathcal J}^\mu{}_{\dot\alpha I}\),
the R–symmetry
currents, and typically a scalar primary at the bottom, subject to
unitarity bounds.

\begin{assumption}[No dynamical gravity (NG)]
\label{ass:4c-NG}
The boundary theory is non–gravitational: it has a unique conserved
stress tensor \(T_{\mu\nu}\) generating translations and does not
contain any additional massless spin–2 gauge field (graviton) or
distinct stress tensors associated with decoupled sectors.
\end{assumption}

In the representation theory of superconformal algebras, one typically
distinguishes between:

\begin{itemize}
  \item \emph{matter} or \emph{Yang–Mills} multiplets, which contain
        fields of spin \(\le1\) (plus the stress tensor in a specific
        short multiplet), and
  \item \emph{supergravity multiplets}, which necessarily include a
        massless spin–2 gauge field (graviton) and additional gauge
        fields fixed by supersymmetry.
\end{itemize}

We use the following standard structural input.

\begin{assumption}[Spin–2 saturation for \(N>4\)]
\label{ass:4c-spin2-saturation}
In four dimensions, any unitary positive–energy representation of
\(\mathfrak{su}(2,2|N)\) with more than \(16\) real supercharges
(i.e.\ \(N>4\)) whose component fields include a unique conserved
stress tensor necessarily coincides with a supergravity multiplet and
thus contains a dynamical graviton.  In other words, for \(N>4\) there
are no purely ``matter'' stress tensor multiplets compatible with
Assumption~\ref{ass:4c-NG}.
\end{assumption}

Assumption~\ref{ass:4c-spin2-saturation} summarises the standard
representation–theoretic lore familiar from the classification of 4D
supergravity multiplets: once the number of supercharges exceeds \(16\),
the stress tensor cannot sit in a purely matter multiplet.

\begin{lemma}[Bound on the number of supercharges]
\label{lem:4c-Nle4}
Under Assumptions~\ref{ass:4c-NG} and
\ref{ass:4c-spin2-saturation}, the superconformal algebra
\(\mathfrak g_{\mathrm{SC}}\) constructed in Subprogram~4b has at most
\(16\) real Poincaré supercharges, i.e.
\begin{equation}
  N\;\le\;4.
\end{equation}
\end{lemma}

\begin{proof}
By Proposition~\ref{prop:4c-su-22N}, \(\mathfrak g_{\mathrm{SC}}\) is a
superconformal algebra of type \(\mathfrak{su}(2,2|N)\) for some
\(N\ge1\), with the
stress tensor \(T_{\mu\nu}\) sitting in a stress tensor multiplet in the
sense of Definition~\ref{def:4c-T-multiplet}.  Suppose, for contradiction,
that \(N>4\).  Then by Assumption~\ref{ass:4c-spin2-saturation}, this
multiplet necessarily contains a dynamical graviton and realises a
supergravity multiplet, contradicting Assumption~\ref{ass:4c-NG} that
our theory is non–gravitational.  Hence \(N\le4\).
\end{proof}

Thus, after imposing non–gravitationality and the standard structure of
superconformal multiplets, the only candidates are \(N=1,2,3,4\).

\subsection{Interacting SK fixed points and exclusion of decoupled free sectors}
\label{subsec:4c-interacting}

We now sharpen the interpretation of \(N\) so that it counts
supersymmetries of the \emph{interacting} universality class singled out
by the SK RG structure, not of arbitrary decoupled free factors.

\begin{assumption}[Interacting SK fixed point]
\label{ass:4c-interacting}
The UV fixed point under consideration is \emph{interacting} in the
following sense:
\begin{enumerate}[(a)]
  \item the SK RG monotone \(\mathcal M\) of Subprogram~3b admits the
        following strengthening at the UV fixed point: it is strictly
        monotone along nontrivial RG flows, additive under decoupled
        tensor products, and has a nondegenerate local extremum there
        (in particular, no flat directions corresponding to decoupled
        generalized free sectors);

  \item the GNS Hilbert space \(\mathcal H\) is generated, as a
        representation of \(\mathfrak g_{\mathrm{SC}}\), by the orbit
        of local operators in the stress tensor multiplet and their
        descendants; in particular, there is no orthogonal direct
        summand carrying a decoupled free superconformal theory with an
        independent stress tensor.

\end{enumerate}
\end{assumption}

In addition we exclude higher–spin enhanced free theories.

\begin{assumption}[No higher–spin enhancement]
\label{ass:4c-noHS}
The UV fixed point carries no infinite tower of conserved higher–spin
currents beyond the stress tensor and its superconformal descendants.
Equivalently, it is not a free or generalized free CFT in the sense of
the higher–spin symmetry no–go theorems.
\end{assumption}

Assumption~\ref{ass:4c-noHS} is motivated by the results showing that
in unitary CFTs in \(d>2\), an infinite tower of higher–spin symmetries
implies that the theory is essentially free.

\begin{lemma}[Intrinsic nature of \(N\) for interacting SK fixed points]
\label{lem:4c-N-intrinsic}
Under Assumptions~\ref{ass:4c-interacting} and
\ref{ass:4c-noHS}, the integer \(N\) appearing in
Proposition~\ref{prop:4c-su-22N} is an intrinsic invariant of the
interacting SK universality class: it cannot be increased by tensoring
with decoupled free superconformal sectors without violating the
standing assumptions.
\end{lemma}

\begin{proof}
Suppose we attempt to pass from our interacting fixed point \(\mathcal C\)
with superconformal symmetry \(\mathfrak{su}(2,2|N)\) to a product theory
\(\mathcal C\otimes\mathcal F\), where \(\mathcal F\) is a decoupled free
SCFT carrying additional supersymmetry \(\mathfrak{su}(2,2|N')\) with
\(N'>0\).  Then:

\begin{enumerate}[(1)]
  \item The product theory has two independent, conserved, symmetric
        stress tensors, \(T_{\mu\nu}^{(\mathcal C)}\) and
        \(T_{\mu\nu}^{(\mathcal F)}\), whose sum generates the total
        translation operator.  This violates the uniqueness of the
        stress tensor in Assumption~\ref{ass:4c-NG}, unless one of the
        factors is trivial.

  \item Alternatively, if one attempts to keep only the diagonal
        combination of the two stress tensors while retaining both sets
        of odd generators as symmetries of the product, the additional
        higher–spin currents in the free factor \(\mathcal F\) violate
        Assumption~\ref{ass:4c-noHS}.

  \item Finally, the SK RG monotone \(\mathcal M\) constructed in
        Subprogram~3b is assumed in
        Assumption~\ref{ass:4c-interacting}(a) to be additive over
        tensor products of decoupled sectors and to have a
        nondegenerate UV extremum; this rules out a nontrivial direct
        product decomposition of the fixed point.
\end{enumerate}

In all cases, raising the effective number of supercharges by adding a
decoupled free factor contradicts the assumptions.  It follows that the
value of \(N\) determined by Proposition~\ref{prop:4c-su-22N} and
Lemma~\ref{lem:4c-Nle4} is an invariant of the interacting SK fixed
point.
\end{proof}

Combined with Lemma~\ref{lem:4c-Nle4}, we thus have:
\begin{equation}
  1 \;\le\; N \;\le\; 4,
\end{equation}
for the interacting, non–gravitational UV fixed point singled out by
our SK construction.

\subsection{Maximality principle and selection of \texorpdfstring{$N=4$}{N=4}}
\label{subsec:4c-maximality}

We now articulate the maximality principle that connects the SK RG
structure to a preference for maximal supersymmetry.

The remainder of this subsection is explicitly conditional. We do not
attempt to derive the maximality principle or the existence statement
below inside the present appendix; we record only what they imply once
the imported superconformal algebra and the previous assumptions are in
place.

\begin{assumption}[Maximality of symmetry at SK fixed points]
\label{prin:4c-maximality}
Consider the class of interacting, non–gravitational, SK–admissible UV
fixed points in four dimensions satisfying the standing Wightman and SK
assumptions and equipped with the two–time BRST/RG structure of
Subprograms~1a and~3b.  Among all such fixed points compatible with a given set
of macroscopic constraints (number of dimensions, spin range, locality,
etc.), the SK RG dynamics selects those with the \emph{maximal} amount
of unbroken symmetry, in the sense that:
\begin{enumerate}[(a)]
  \item the SK RG monotone \(\mathcal M\) attains local extrema on
        fixed points whose symmetry algebra is maximal within their
        basin of attraction in theory space, and

  \item for fixed macroscopic data, increasing the number of
        supercharges \(N\) (within the allowed algebraic range) cannot
        be compensated by adding further continuous symmetries without
        changing the universality class.
\end{enumerate}
\end{assumption}

Assumption~\ref{prin:4c-maximality} is an additional assumption of this
optional strengthening. It promotes the observed tendency of RG flows to land on
highly symmetric fixed points, together with the SK construction of odd
symmetries from monotones, into a symmetry-selection rule. It is not
derived here.

\begin{assumption}[Existence of interacting SCFTs for all \(N\le4\)]
\label{ass:4c-existence}
For each \(N\in\{1,2,3,4\}\) there exist interacting, non–gravitational,
unitary four–dimensional superconformal field theories with symmetry
\(\mathfrak{su}(2,2|N)\) satisfying the standing assumptions (in
particular, with a unique stress tensor and compact R–symmetry).
\end{assumption}

Assumption~\ref{ass:4c-existence} is motivated by the known families of
$\mathcal N=1,2,4$ SCFTs and the recent constructions of genuinely
$\mathcal N=3$ theories; we do not require a complete classification,
only that the possible $N$ in the allowed range are all realised.

\begin{proposition}[Maximal supersymmetry: \(N=4\)]
\label{prop:4c-N-equals-4}
Let \(\mathcal C\) be an interacting, non–gravitational, SK–admissible
UV fixed point in four dimensions, equipped with the two-time RG
structure and SK monotone of Subprogram~3b, and suppose
Assumptions~\ref{ass:4c-NG}, \ref{ass:4c-spin2-saturation},
\ref{ass:4c-interacting}, \ref{ass:4c-noHS},
\ref{ass:4c-existence}, and
Assumption~\ref{prin:4c-maximality} hold.
Then the superconformal symmetry of \(\mathcal C\) is maximally
supersymmetric:
\begin{equation}
  \mathfrak g_{\mathrm{SC}} \;\cong\; \mathfrak{su}(2,2|4),
\end{equation}
i.e.\ the boundary theory has \(\mathcal N=4\) superconformal symmetry.
\end{proposition}

\begin{proof}
By Proposition~\ref{prop:4c-su-22N} and Lemma~\ref{lem:4c-Nle4}, the
superconformal symmetry is of imported type \(\mathfrak{su}(2,2|N)\) with
\(1\le N\le4\).  By Lemma~\ref{lem:4c-N-intrinsic}, \(N\) is an intrinsic
invariant of the interacting SK universality class.  By
Assumption~\ref{ass:4c-existence}, there exist interacting SCFTs
realising each \(N\) in this range; in particular, there exist candidate
fixed points with superconformal symmetry \(\mathfrak{su}(2,2|4)\).

Assumption~\ref{prin:4c-maximality} then asserts that, within the class
of SK–admissible fixed points compatible with a given set of macroscopic
constraints, the RG dynamics selects fixed points with maximal symmetry.
Specialising to the superconformal symmetries under discussion, this
implies that the emergent algebra must correspond to the largest
available value of \(N\), namely \(N=4\).  Hence
\(\mathfrak g_{\mathrm{SC}}\cong\mathfrak{su}(2,2|4)\).
\end{proof}

\begin{remark}[Projective form and R–symmetry]
\label{rmk:4c-psu}
The even part of \(\mathfrak{su}(2,2|4)\) is
\(\mathfrak{so}(4,2)\oplus\mathfrak{su}(4)\oplus\mathfrak u(1)\).  In
many physical applications the central \(\mathfrak u(1)\) acts trivially
and one works with the projective algebra
\(\mathfrak{psu}(2,2|4)\).  Our arguments are insensitive to this subtle
distinction, which can be fixed later by imposing the absence of
conserved currents for the central \(\mathfrak u(1)\).
\end{remark}

\subsection{Summary of Subprogram 4c}
\label{subsec:4c-summary}

We summarise the logical content of this subprogram.

\begin{itemize}
  \item We import from Subprogram~4b the same-\(\mathcal H\)
        realization of a \(4\)D superconformal algebra, written in
        standard form as \(\mathfrak{su}(2,2|N)\) (with the projective
        quotient in the \(N=4\) case)
        (Proposition~\ref{prop:4c-su-22N}).

  \item Under the additional non--gravitationality and spin--2
        saturation assumptions, the number of supercharges is bounded by
        \(N\le4\)
        (Lemma~\ref{lem:4c-Nle4}).

  \item Under the additional interaction, nonfactorization, and
        no--higher--spin assumptions, \(N\) is an intrinsic invariant of
        the interacting SK universality class and cannot be raised by
        tensoring with decoupled free factors
        (Lemma~\ref{lem:4c-N-intrinsic}).

  \item Finally, under the additional maximality principle and the
        existence of interacting SCFTs for each \(N\le4\), the emergent
        symmetry at the UV fixed point must be maximally
        supersymmetric,
        \(\mathfrak{su}(2,2|4)\)
        (Proposition~\ref{prop:4c-N-equals-4}).

\end{itemize}

This subsection therefore gives a conditional route from the imported
superconformal closure of Subprogram~4b to maximal
\(\mathcal N=4\) symmetry under the added assumptions above.
The conjecture below records the corresponding cosmology-theory
interpretation suggested by this optional strengthening.
\begin{conjecture}[String realisation of a \(\mathbf{2T}_{\mathrm{Main}}\) cosmology theory]
\label{conj:4c-string-realisation}
A cosmology theory that implements the \(2T\)-Main interface
\(\mathbf{2T}_{\mathrm{Main}}\) of Definition~\ref{def:2T-Main} is realised by
the \(\mathcal N=4\) superconformal field theory with symmetry algebra
\(\mathfrak{su}(2,2|4)\), together with its type IIB superstring dual on
\(\mathrm{AdS}_5\times S^5\).
\end{conjecture}
